\begin{frame}{그림으로: $x=1$과 $x=5$의 우도 차이}
  \begin{center}
  \includegraphics[height=5.6cm]{gen_likelihood_x1_x5.png}
  \end{center}
  \vskip0.3em
  \small
  진짜 $p(x) = \mathcal{N}(x;0,2)$ 위에서 $x=1$(두 봉우리 중 하나에
  가까움)과 $x=5$(어느 봉우리에도 속하지 않음)의 높이를 직접 비교 ---
  $x=5$ 쪽이 \textbf{훨씬 낮다}. 근사분포 $q(z|x)$가 이만큼 다른
  $x$들을 전부 감당하려면 \textbf{$x$마다 다른 분포}가 필요.
\end{frame}
